\section{Dense representation and stationary Markov realization}

\subsection{Representation invariance}
Let $V\in\mathbb R^{D\times D}$ be invertible with rows $v_j^\top$. The
observed covariate is $X=Sv_J$, the latent label is $S\theta_J$, and the ambient
teacher is $w^\star=V^{-1}\theta$.

\begin{proposition}
For $\widehat w=V^{-1}\widehat\theta$,
\[
\E\langle X,\widehat w-w^\star\rangle^2
=\sum_{j=1}^D\lambda_j(\widehat\theta_j-\theta_j)^2.
\]
If $\lambda_1\ge\cdots\ge\lambda_D$ and
$mI\preceq VV^\top\preceq LI$, then for
$\Sigma=V^\top\Lambda V$ and $\Lambda=\operatorname{diag}(\lambda)$,
\[
m\lambda_j\le\lambda_j(\Sigma)\le L\lambda_j
\quad (j=1,\ldots,D),
\]
where eigenvalues are ordered decreasingly.
\end{proposition}

\begin{proof}
Since $v_j^\top\widehat w=(V\widehat w)_j=\widehat\theta_j$, the prediction
identity is immediate. The nonzero eigenvalues of $V^\top\Lambda V$ equal
those of $\Lambda^{1/2}VV^\top\Lambda^{1/2}$. The assumed operator inequality
implies
\[
m\Lambda\preceq\Lambda^{1/2}VV^\top\Lambda^{1/2}\preceq L\Lambda.
\]
Eigenvalue monotonicity under Loewner order gives the claim.
\end{proof}

\subsection{Why dense autoregression is a negative control}
Consider the finite-dimensional Gaussian process
\[
x_{t,j}=\rho_jx_{t-1,j}
+\sqrt{1-\rho_j^2}\sqrt{\lambda_j}z_{t,j},
\qquad |\rho_j|<1,
\]
with independent standard Gaussian innovations.

\begin{proposition}[Finite-dimensional exact recovery]
For every $N\ge D$, the $N\times D$ design matrix has full column rank almost
surely. Consequently, noiseless least squares recovers every teacher exactly,
regardless of the mode-dependent autocorrelation times.
\end{proposition}

\begin{proof}
The stacked entries of the design matrix have a nonsingular joint Gaussian
law: the map from the stationary initial state and subsequent innovations to
the trajectory is invertible, and every innovation variance is positive.
Hence the design has a density with respect to Lebesgue measure on
$\mathbb R^{N\times D}$. Rank-deficient matrices form an algebraic set of
Lebesgue measure zero when $N\ge D$. Full column rank makes the noiseless least
squares solution unique and equal to the teacher.
\end{proof}

This proposition does not describe an infinite-dimensional joint limit. It
serves only to show that mode-wise correlation, without restricted innovation
coverage, cannot universally force the renewal exponent.

\subsection{A reversible nonseparable Markov chain}
Consider states $(j,s)$, $s\in\{-1,1\}$, with transition kernel
\[
P((j,s),(k,t))=
(1-\tau_j^{-1})\1\{(j,s)=(k,t)\}
+\tau_j^{-1}\frac{q_k}{2}.
\]

\begin{proposition}
The chain is reversible with stationary law $\pi_{j,s}=\lambda_j/2$. For
$X_t=S_te_{J_t}$,
\[
\E[X_tX_{t+\ell}^\top]
=\sum_j\lambda_j(1-\tau_j^{-1})^\ell e_je_j^\top.
\]
Unless all $\tau_j$ are equal, this covariance is not a scalar time kernel
times the spatial covariance.
\end{proposition}

\begin{proof}
For distinct states,
\[
\pi_{j,s}P((j,s),(k,t))
=\frac{\lambda_jq_k}{4\tau_j}
=\frac{Zq_jq_k}{4},
\]
which is symmetric; diagonal transitions are immediate. Conditional on
$J_t=j$, the only contribution correlated with $X_t$ is the event of no
refresh during $\ell$ steps, with probability
$(1-\tau_j^{-1})^\ell$. Any refresh resamples a mean-zero sign.
\end{proof}

\subsection{Matched global autocorrelation with different exponents}
For the reversible chain above, define the trace integrated autocorrelation
\[
\tau_{\rm tr}
=1+2\sum_{\ell\ge1}
\frac{\operatorname{tr}\E[X_tX_{t+\ell}^\top]}
{\operatorname{tr}\E[X_tX_t^\top]}.
\]
Since $\sum_j\lambda_j=1$, the covariance formula and the geometric series give
\begin{align*}
\tau_{\rm tr}
&=1+2\sum_j\lambda_j\sum_{\ell\ge1}(1-\tau_j^{-1})^\ell\\
&=1+2\sum_j\lambda_j(\tau_j-1)
=2\sum_j\lambda_j\tau_j-1.
\end{align*}

Assume $\lambda_j=j^{-a}/\zeta(a)$ and choose $0<r<a-1$.  Then
\[
m_r:=\sum_j\lambda_jj^r=\frac{\zeta(a-r)}{\zeta(a)}<\infty.
\]
For $T\ge m_r$, both $\tau_j^{\rm unif}=T$ and
$\tau_j^{\rm align}=(T/m_r)j^r$ are at least one, and both give
$\tau_{\rm tr}=2T-1$.  On the other hand,
\[
q_j^{\rm unif}
=\frac{\lambda_j/T}{\sum_k\lambda_k/T}
=\frac{j^{-a}}{\zeta(a)},
\]
whereas
\[
q_j^{\rm align}
=\frac{j^{-(a+r)}}{\zeta(a+r)}.
\]
The finite-block exponents follow from the power-law theorem.  For stationary
raw horizons, the aligned residual exponent is $(a-1)/r>1$.  Under
$1<b<a+r+1$, the innovation exponent $(b-1)/(a+r)$ is below one, so the
residual term decays faster and the same two distinct slopes remain.  This
proves the matched-autocorrelation proposition in the main paper.

\subsection{Compute-optimal allocation}
Let $d=b-1>0$ and $p=a+r$.  The noiseless risk is comparable to
$M^{-d}+B^{-d/p}$.  Under $MB\le C$, choose
$M\asymp C^{1/(p+1)}$ and $B\asymp C^{p/(p+1)}$ to obtain
$O(C^{-d/(p+1)})$.  Conversely, for any feasible pair, at least one of
$M\le C^{1/(p+1)}$ or $B\le C^{p/(p+1)}$ holds.  In the first case
$M^{-d}\ge C^{-d/(p+1)}$; in the second,
$B^{-d/p}\ge C^{-d/(p+1)}$.  Hence the rate and both allocation exponents are
sharp up to constants.
